\section{Projected tables and the second-moment reduction}
\label{sec:tables}

We now turn the selected rows and their birth prescriptions into weighted
interval tests.  A row initially appears on a finer grid than its true
numerator grid.  As the remaining primes are read, we restrict that finer
grid and adjust its weights.  The adjustment preserves first mass and
cancels many close-pair fluctuations.  The uncancelled fluctuations are
controlled by squared masses of classes of labels.

The arithmetic estimates needed for this procedure are stated precisely
below as Propositions~\ref{tab:first-join-input}--\ref{tab:partition-input}.
Their proofs occupy Sections~\ref{sec:direct}--\ref{sec:joins}.  The present
section defines the procedure and proves the second-moment reduction
conditional on those inputs.  In particular, none of their counting
conclusions is presumed to follow merely from the table notation.

\subsection{Chronology, labels, and the translation core}

Activate a centre $L_i$ at its start height $Y_0(i)$.  At the beginning of
each dyadic height $Y$, join active centres $L_i,L_j$ by an edge when
$\log UT\le\beta Y$, where, as before,
\[
 G=(L_i,L_j),\qquad U=L_i/G,\qquad T=L_j/G.
\]
The resulting graph components grow and coalesce as $Y$ increases.  All
prime powers in a component at primes with $\log\log p\ge Y$ are equal:
an exponent difference at such a prime would already exceed the edge
threshold.  Equality propagates along paths.

Every active band satisfies $B_i\le Y/A_0$.  The centre count from
Section~\ref{sec:centres} therefore bounds the number of vertices by
$\exp(4Y/A_0+O(1))$.  Sum edge log-ratios along a path and use
$|\log k_{L_i}|=O(B_i)$ to obtain, for any two vertices of one component,
\begin{equation}
 \log\bigl(2+UT+K+r_{\max}/r_{\min}\bigr)
 \le \exp(h_*Y/10).
 \label{eq:component-path}
\end{equation}
Here $K=[L_i,L_j]r_{\max}$, and the parameter hierarchy makes $A_0$
sufficiently large for this bound.  Read the common tail primes of the
bin in increasing order.  The block is finite, so finitely many height
steps, including empty ones if necessary, process every prime and start
and eventually join every centre.

At an intermediate time let $I$ be the common remaining prime-power
part of a component.  A row assigned to $L_i$ has a processed deficit
$d$ and a remaining deficit $a\mid I$, with
\[
 v=\frac{L_i}{da}.
\]
Its current projected grid is
\begin{equation}
 x=\frac{z+a\gamma}{L_i/d},\qquad z\in\mathbb Z.
 \label{tab:projected-grid}
\end{equation}
We distinguish labels even when their locations coincide.  A label
records the centre, the assigned row, and the numerator coordinate.
Each label has a raw weight $w_x\in[0,1]$.  Its weighted point mass is
\[
 \lambda_x=\pi_I(a)w_x.
\]
We also write $\lambda$ for the atomic point measure with these
label masses; coincident labels retain their multiplicities.
All point sums are per unit length.  The first mass of an array is
$\sum_x r_x\lambda_x$, with $r_x=r_i$ on a row of centre $i$.
We call the collection of these labelled arrays and raw weights for
one component its \emph{projected table}.

The role of the mask law is expressed by the exact identity
\begin{equation}
 r_i\pi_I(a)\frac{L_i}{d}
 R\left(\frac{L_i/I}{d}\right)
 =r_i\phi(v)=k_{L_i}\pi_{L_i}(da).
 \label{tab:first-currency}
\end{equation}
Thus the single coefficient $\pi_I(a)$ in first mass is distinct from
a numerator-density factor.  We sum only over the assigned good masks;
we do not interpret their mask coefficients as automatic numerator
exclusions.

At birth the raw weights are the static prescriptions of
Section~\ref{sec:prescriptions}, including the additional static test
constructed in Section~\ref{sec:direct}.  Subsequent locations descend
from labels of the same row at the same physical point.  For example,
if a further mask portion $b$ is processed, the old denominator, phase,
and numerator are all $b$ times their new values.  Raw weights never
increase.  Dead labels may be omitted or assigned weight zero.  We use
the circle for partitions and the periodic lifts for pair counts;
these give the same nearby pairs because all physical widths are small.

The order at a height is as follows: insert newborn projections,
perform the partition tests when a component is created or changes,
clip the resulting cell loads if needed, and then perform its
sparse-prime updates.  We call the birth factors \emph{static}; the
prime updates from its start onward are \emph{adaptive}.

For the translations used below we need a common symmetry group.
Immediately before a prime $p$ in bin $Y$, put
\[
 X=\exp(2h_HY).
\]
Start with the gcd of the centres in the component.  At each already
passed prime $l$, including all primes in earlier bins, reduce its
exponent by
\[
 4+\left\lfloor\frac{CX}{\log l}\right\rfloor,
\]
clipping at zero, where $C$ is a sufficiently large fixed constant.
The resulting integer is the \emph{pre-$p$ core}.  After the update,
make this reduction at $p$ as well; the result is the
\emph{post-$p$ core}, denoted $\mathrm{core}_p$.  At bin start only
earlier primes have been reduced.  Cores decrease by divisibility
through this chronology, including at joins.

The structural requirement on every operation is invariance, or
equivariance of its labelled objects, under translation by the
reciprocal of the current core.  The birth prescriptions have this
property.  Indeed, a processed deficit uses at most
$1+\lfloor C_{\rm mask}B_i/\log l\rfloor$ powers at $l$, a hard
exclusion costs one residue digit, and a soft condition need only know
the numerator modulo
\[
 l^{\min\{\nu_l(V),\,1+\lfloor h_s\log H_v/\log l\rfloor\}}.
\]
Here $\log H_v\ll X$ for every model in use.  The extra static tests
use squarefree low numerator moduli.  These observations show both
that the core divides the projected denominator and that reciprocal
core translation preserves the birth data.  The construction of
partitions and the matching rule below preserve this symmetry.

The loss of the core is bounded by
\begin{equation}
 \log\frac{L_i}{\mathrm{core}_p}
 \ll \exp(h_*Y/10)+p+\exp(O(X))\ll p.
 \label{eq:core-loss}
\end{equation}
For the gcd loss, sum the edge log-ratios over a spanning tree.  At a
fixed prime the maximal exponent loss needed to reach the common gcd
is at most the sum of the absolute exponent changes on that tree.
The remaining loss follows from the prime-log bound and the reductions
just specified.  For instance, primes not exceeding $\exp(O(X))$
give the last term; above that threshold the sum of the terms
$4\log l+CX$ up to $p$ is $O(p)$.  Since $\log p\ge\exp(Y)$ and
$2h_H<1$, the last inequality follows.  The path term is bounded using
\eqref{eq:component-path} and the choice of $A_0$.

In particular, if $p^\nu\parallel L_i$ is a currently unprocessed
tail prime power, then $p^\nu$ divides the pre-$p$ core.  All previous
decisions therefore preserve translation by $1/p^\nu$.  This includes
the transport of earlier projection data along the same-location
ancestor of a current label: at every earlier time that prime power
belonged to the core of its then-component.

\subsection{Cells, clipping, and a prime matching rule}

The construction uses, separately for each exact width in a component,
a partition of the circle into cells.  A partition is replaced only
at creation or change of that component.  The precise existence and
loss estimate for these partitions is
Proposition~\ref{tab:partition-input}.  We first state the properties
needed to define the update.
At a change, the previous components contributing existing tables are
called the \emph{old children}; newly inserted centres are not old
children.

After the partition tests, any surviving points in the same cell,
across all masks, are separated by
\begin{equation}
 O\bigl(r_i e^{-E_gB_i}\bigr)
 \label{tab:cell-span}
\end{equation}
when measured at an entry of centre $i$ there.  For a high-type
component having direct newborn comparisons at height $Y$, the
stronger bound $O(r_i e^{-E_gY})$ holds.  Cells are padded at surviving
points: in a later sparse step at height $Y'$ of the same unchanged
component, same-width pairs at distance at most
\[
 r_i\exp\left(-\tfrac12\exp(\sigma_gY')\right)
\]
lie in the same cell.  At a component change, survivors from one old
child that lie in one new cell come from one of that child's old
cells.

For each current remaining mask separately, cap the total raw load
in a cell at $N_0$ by scaling all its weights if necessary.  If $x_0$
is its pre-clip load and $x_D$ is the contribution of an old child $D$,
then the inherited cap gives
\begin{equation}
 \sum_Dx_D^2\le N_0x_0.
 \label{tab:clip-child}
\end{equation}
When $x_0>N_0$, it follows that
\[
 x_0-N_0
 \le\frac{x_0^2-\sum_Dx_D^2}{N_0}.
\]
Consequently the first mass lost in clipping is at most $N_0^{-1}$
times the same-width, same-mask pair count in the new cells, excluding
pairs from the same old child and including newborn diagonals.  Its
currency is
\begin{equation}
 r_i\pi_I(a)w_xw_y,
 \label{tab:one-law}
\end{equation}
with \emph{one} mask-law factor.  This is the narrow count in
Proposition~\ref{tab:first-join-input}.

There is also an equivalence relation on surviving labels, separately
for each exact width but not separately for each mask.  Its
equivalence classes, called \emph{accounting classes}, lie inside
cells.  At a component change retain each old-child class restricted
to its surviving labels; newborn labels start as singletons.  At
each sparse-prime step coarsen these classes only after updating
weights, using the tightly aligned successful plus pairs selected in
Proposition~\ref{tab:sparse-input}.  At every thinning or projection
we restrict a class to its surviving descendants without splitting
its equivalence relation further.  In particular, deleting an old
connecting label does not split a class.  The new partitions must
respect all these inherited equivalences.

Fix a tail prime power $p^\nu\parallel L_i$.  Good masks have no
deficit of exponent at least two at this prime.  For each mask after
$p$, and in each cell of a fixed width, form two lists before the
update.  Write $I=p^\nu I'$ and fix an after-prime mask $a\mid I'$.
The \emph{plus list} has before-prime mask $pa$, with the single
omission at $p$; the \emph{minus list} has before-prime mask $a$,
with the full factor.  First multiply plus weights by $R_p$ if $\nu\ge2$.
A plus label succeeds exactly when its before-prime numerator is
divisible by $p$, so that it lies on its next projection.

Fractionally match the residual-scaled plus weights to the minus
weights in that cell.  Concretely, use a finite nonnegative matrix
whose row and column sums are bounded by the endpoint weights and
whose total mass is the smaller of the two list loads.  Product
coupling provides such a matrix equivariantly.  Retain a matched
plus portion precisely on success, and retain its matched minus
portion precisely on failure of that plus label.  Retain unpaired
plus mass on success, and multiply unpaired minus mass by $R_p$.
This rule never moves a retained location and is post-core
equivariant.

\begin{lemma}[One-prime mass and cap preservation]
\label{tab:prime-update}
The matching rule preserves the cap on raw cell load per after-prime
mask.  On averaging over pre-core $p$-translations, each transported
label's weighted mass after the update has mean equal to its
before-prime weighted mass.  In particular, the total first mass per
unit length is preserved by every sparse update.
\end{lemma}

\begin{proof}
Let $P$ be the residual-scaled plus load and $M$ the minus load.
Matched portions contribute their full matched mass between the two
lists.  Only the larger list has unpaired mass, which can only
decrease.  The post-update load is therefore at most $\max(P,M)$,
and hence at most $N_0$.

For these fixed before-prime and after-prime masks, the probabilities
satisfy
\begin{equation}
 \pi_I(a)=R_p\pi_{I'}(a),\qquad
 \pi_I(pa)=\frac{R_p^{[\nu\ge2]}}p\pi_{I'}(a).
 \label{tab:mask-transition}
\end{equation}
Translation by $1/p^\nu$ changes the plus numerator by a unit
modulo $p$.  Its success indicator $S$ thus has orbit mean $1/p$.
The pre-link weights and the matching matrix transport equivariantly
on this orbit.  A matched portion of residual-scaled size $m$ has,
in units of $\pi_{I'}(a)$, before-prime masses $m/p$ and $mR_p$ and
after-prime masses $mS$ and $m(1-S)$.  Unpaired portions give the
same marginal identities.  Summing over the finite translation
orbits proves first-mass preservation.  The orbit is a simultaneous
translation of all lists; the assertion does not require independent
success colors.
\end{proof}

\subsection{The three quantitative inputs}

We specify the counts with their different pair-weight conventions
explicit.  For labels $x,y$ of widths $r_x,r_y$, write
$r_{xy}^- =\min(r_x,r_y)$ and $r_{xy}^+=\max(r_x,r_y)$.
The broad pair currency is
\begin{equation}
 r_{xy}^-\lambda_x\lambda_y
 =r_{xy}^-\pi_I(a_x)\pi_I(a_y)w_xw_y.
 \label{tab:product-law}
\end{equation}
The product of probabilities is used even when $a_x=a_y$.  A centre
pair has a \emph{first-sharing time} when its two centres first lie
in one component; a self-pair has its birth as first-sharing time.
At a join, counts involving new labels are taken before the fresh
cap clipping and after its preliminary partition tests.  For the
broad majorants those tests can be dropped.  For the narrow
majorants they first supply the cell-span restriction, after
which the point weights can be bounded by their incoming values.

The distinction between these currencies is essential. Broad energy
compares two weighted arrays and therefore multiplies their two mask
coefficients. Clipping controls raw load at one fixed mask, so it
uses that mask coefficient once. Table~\ref{tab:currency-map} records
these conventions for reference throughout the remaining proof.

\begin{table}[htbp]
\centering
\renewcommand{\arraystretch}{1.2}
\begin{tabular}{@{}p{.25\linewidth}p{.40\linewidth}p{.27\linewidth}@{}}
\toprule
Quantity & Weight per entry or ordered pair & Restriction \\
\midrule
First mass & $r_x\pi_I(a_x)w_x$ & One label \\
Broad pair count & $r_{xy}^-\pi_I(a_x)\pi_I(a_y)w_xw_y$
  & Both mask factors, even if $a_x=a_y$ \\
Clipping pair count & $r\pi_I(a)w_xw_y$
  & $r_x=r_y=r$, $a_x=a_y=a$ \\
\bottomrule
\end{tabular}
\caption{Mass conventions. Here $I$ is the current common remaining
prime-power part, $a_x$ is the remaining deficit of label $x$, and
$w_x$ is its raw weight. The clipping estimate also uses the new-cell
and old-child restrictions stated below.}
\label{tab:currency-map}
\end{table}

\begin{proposition}[First-sharing comparison input]
\label{tab:first-join-input}
For the selected finite blocks, there is a static additional birth
test, depending only on the row data and squarefree low numerator
moduli and not on the running tables, with arbitrarily small relative
first-mass loss against the hard birth means, such that the following
statements hold.

For any fixed required absolute $C$, the sum, over first-sharing
times and centre pairs, of the broad currencies
\eqref{tab:product-law} for point pairs at distance at most
$Cr_{xy}^+$ is $O(1)$.

For exact same-width, same-remaining-mask point pairs in one new
cell, sum the currencies \eqref{tab:one-law}, excluding pairs from
the same old child and including all newborn self and diagonal
contributions.  This sum is
\begin{equation}
 \ll W+O_{k_0}(W^{1+c})
 \label{tab:narrow-input}
\end{equation}
for a fixed $c>0$.  Both bounds are uniform in the finite block and
in prior table histories obeying the stated operations and core
symmetries.  The constants in \eqref{tab:narrow-input} may depend on
the preceding fixed cutoffs but not on $N_0$.
\end{proposition}

The direct adjacent simultaneous-birth part, including the static additional
test, is proved in Proposition~\ref{dir:birth-comparisons}.  The
remaining first-sharing times are treated in
Section~\ref{sec:joins}.  For the narrow estimates these proofs use
the span bounds \eqref{tab:cell-span}, rather than any unspecified
independence of cell membership.

\begin{proposition}[Sparse-alignment input]
\label{tab:sparse-input}
For each sparse-prime step in bin $Y$, there is an equivariant
selection of exceptional same-width successful plus pairs at
distance at most
\[
 r_i\exp\left(-\tfrac12\exp(\sigma_gY)\right).
\]
Coarsen accounting classes by every nonexceptional pair in this
range whose two arriving masses are retained.  Uniformly under the
allowed prior histories, the following three sums over components,
steps, applicable centre pairs, masks, and locations are bounded.
\begin{enumerate}
\item Actual simultaneous plus successes at distances between
$r_{xy}^+\exp(-\exp(\sigma_gY))$ and $Cr_{xy}^+$, weighted by
$r_{xy}^-$ times the product of their raw arrival weights and their
two after-prime mask probabilities.
\item Before-prime plus pairs at distance at most $Cr_{xy}^+$,
weighted by $r_{xy}^-/p^2$ times the product of their before-prime
raw weights and their two after-prime mask probabilities.
\item The selected exceptional actual successful pairs in the
relaxed same-width range above, with the same product-law currency
as in the first item.
\end{enumerate}
The constants are uniform in the block, in $N_0$, and in the
allowed preceding histories.  In the first and third items one may
use raw arrival weights immediately after the update and before
coarsening.  The successful linked and unpaired portions of a plus
entry have total mass at most that raw arrival weight.  The first
two items impose no equality of widths, and $C$ may be any required
fixed absolute constant.
\end{proposition}

Section~\ref{sec:sparse} gives an explicit selection by deterministic
trains of short intervals and proves these bounds.  Thus the
exception selection is part of the construction, not a choice made
after seeing the total energy.

\begin{proposition}[Chronological partition input]
\label{tab:partition-input}
The cells and accounting classes described above can be constructed
chronologically for every finite block.  The construction preserves
current-core equivariance, the cell spans and padding stated in
\eqref{tab:cell-span} and the paragraph following it, all inherited
accounting classes without splitting, and the property that one
new cell meets at most one surviving old cell of each old child.
After clipping, the raw cap per cell and remaining mask is $N_0$.
It supports the coarsening rule of
Proposition~\ref{tab:sparse-input} at each subsequent sparse step.

Aside from the cap clipping separately bounded by
\eqref{tab:narrow-input}, its total first-mass loss can be made
arbitrarily small relative to the selected block mass by the fixed
parameter cutoffs.  An expected loss bound over the specified
finite partition jitters suffices; the structural assertions and
core equivariance hold for every realized construction.  The
comparison bounds apply to the actual running tables or to the
nonnegative majorants indicated in their proofs.
\end{proposition}

This proposition is proved in Section~\ref{sec:partitions}.  Its
construction only uses previously generated interval trains and
inherited classes.  Conversely, the sparse estimates apply to the
history before the current coarsening.  The two inputs are therefore
used by finite chronological induction: the partition at a change
is constructed from the past, and the next sparse estimate is
applied before its new classes are joined.  No future partition is
used to establish a past estimate.

More explicitly, the order
is inherited tables and newborns, preliminary history tests,
realized padded cells and their further tests, cap clipping,
prime matching updates, and then class coarsening.

The history
tests use deterministic supplies of past interval trains and
single-row upper fields, not the cap's value.  The conditional
jitter bounds at each height charge a proportion of entering
first mass, which at that height is at most $W$; their
height-dependent failure probabilities are summable.  The
first-sharing and sparse
comparison bounds use weights at most one and their indicated
prefix or upper-field structure.  Thus their constants and the
non-clipping loss budgets are fixed before $N_0$ is chosen.

Figure~\ref{fig:chronology} separates the operations from the estimates
that justify them. Its return to previously stored history is the reason
that the partition and sparse-step arguments can be used in one finite
induction.

\begin{figure}[htbp]
\centering
\newcommand{\chronobox}[1]{\fbox{\parbox{.88\linewidth}{\centering #1}}}
\chronobox{\textbf{Component creation or change}\\
Old tables, accounting classes and stored trains; fixed newborn arrays}
\\[3pt] $\downarrow$ \\[3pt]
\chronobox{History tests and padded cells preserving old classes\\
Section~\ref{sec:partitions}}
\\[3pt] $\downarrow$ \\[3pt]
\chronobox{Cap clipping\\
First-sharing estimates in Sections~\ref{sec:direct} and~\ref{sec:joins}
 bound the lost first mass}
\\[3pt] $\downarrow$ \\[3pt]
\chronobox{\textbf{Next prime in the component}\\
Match plus and minus portions; perform the divisibility update\\
Lemma~\ref{tab:prime-update}}
\\[3pt] $\downarrow$ \\[3pt]
\chronobox{Count arrivals; store the full prescribed train supply\\
Section~\ref{sec:sparse}}
\\[3pt] $\downarrow$ \\[3pt]
\chronobox{Coarsen classes at accepted close pairs\\
Continue to the next prime, or use this stored history at the next change}
\caption{Chronology of the finite construction. A partition uses only
past train intervals; a sparse-step estimate uses the partition already
in place. The comparison estimates bound the cost of the displayed
operations and do not change the arrays.}
\label{fig:chronology}
\end{figure}


\subsection{Accounting energy and broad pair energy}

We now show what these inputs accomplish.  Write
\[
 \mathcal B
 =\sum_{\text{current components}}
   \sum_{\substack{x,y\text{ in that component}\\
               \operatorname{dist}(x,y)\le10r_{xy}^+}}
       r_{xy}^-\lambda_x\lambda_y
\]
for the broad energy.  Discards decrease it.  Births and first joins
inject the pairs counted in
Proposition~\ref{tab:first-join-input}.  We must bound the remaining
increments from sparse updates.

For this purpose define the accounting energy
\begin{equation}
 \mathcal A=\sum_A r_A\lambda(A)^2,
 \qquad \lambda(A)=\sum_{x\in A}\lambda_x,
 \label{tab:accounting-energy}
\end{equation}
where $A$ ranges over current accounting classes, all of whose
labels have the same width $r_A$.  Since a class is inside one cell,
the cap for each mask and $\sum_{a\mid I}\pi_I(a)=1$ give
$\lambda(A)\le N_0$.  This applies to inherited old classes even
before a fresh clipping.

\begin{lemma}[Variance budget]
\label{tab:variance-budget}
Let $J_p$ be the increment of \eqref{tab:accounting-energy} on old
classes during a sparse-prime update, before coarsening, and let
$K_p'\ge0$ be the increment from the subsequent coarsening.  Then
\[
 J_p\ge0,\qquad
 \sum_p(J_p+K_p')\ll N_0W,
\]
where the sum includes all component steps.
\end{lemma}

\begin{proof}
Transport old labels and classes by the pre-core translations
$1/p^\nu$, assigning zero post-weight to vanished labels.
For a residual-scaled matched portion let $\mu$ denote its raw
mass times its after-mask probability.  Its contribution to the
increment of a class $A$ is
\[
 \mu\bigl(\mathbf1_{\{\text{plus}\in A\}}
          -\mathbf1_{\{\text{minus}\in A\}}\bigr)(S-1/p).
\]
For an unpaired plus portion omit the minus indicator; an unpaired
minus portion has zero increment.  This formula incorporates the
factor $R_p^{[\nu\ge2]}$ already absorbed into the residual-scaled
mass, and applies separately with each portion's own mask probability.
Thus it also applies to classes mixing masks and to links crossing
classes.

Let $\mathcal G$ be the finite group generated by $1/p^\nu$.
Write $Y_A(g)$ for the post-mass on the transported class $gA$,
and $m_A$ for its old mass, which is constant on its class orbit.
The formula above gives $\mathbb E_{g\in\mathcal G}Y_A(g)=m_A$.
Since each $g$ permutes the complete set of old classes per unit
length,
\[
 J_p=\sum_A r_A\bigl(\mathbb E_gY_A(g)^2-m_A^2\bigr)
     =\sum_A r_A\operatorname{Var}_g Y_A(g)\ge0.
\]
Possible stabilizers cause no change: averaging over the whole group
repeats every element of a class orbit equally.  This is an identity
for the actual summed energy increment, not an extra randomized
construction.  Coarsening adds the cross products between distinct
old classes that are joined, giving $K_p'\ge0$.

The total first mass inserted at births is at most $W$, by the
hard mean upper bound and \eqref{tab:first-currency}.
Lemma~\ref{tab:prime-update} preserves first mass at sparse steps,
so the total first mass lost in all discards is at most $W$.
Discarding class mass $d_A$ drops its square by at most
$2N_0d_A$.  Restriction to survivors creates no further drop through
splitting, because the equivalence relation is retained.  New
singleton classes contribute nonnegative injections, and the final
accounting energy is at most $N_0W$.  Telescoping the energy gives
the asserted budget.
\end{proof}

\begin{proposition}[Broad-energy bound]
\label{tab:energy-bound}
Assuming Propositions~\ref{tab:first-join-input}--\ref{tab:partition-input},
the final broad energy is bounded uniformly down the finite-block
tail.  The bound may depend on the fixed cap $N_0$ and preceding
hierarchy parameters, but not on the block.
\end{proposition}

\begin{proof}
\emph{Orbit expansion.}\quad
Index the centered fluctuations at a prime by residual-scaled plus
portions.  A portion of raw size $m$ and after-mask probability
$\pi$ has centered coefficient
\[
 \pi m(S-1/p).
\]
For an unpaired portion it acts only at its plus endpoint.  For a
matched portion it acts there and with the opposite sign at its
minus endpoint of the same width.  Indeed, in units of $\pi$, the
matched pre-masses are $m/p,m(1-1/p)$ and the post-masses are
$mS,m(1-S)$.

The broad distance kernel is invariant under simultaneous
translation.  On expanding the bilinear energy and averaging over
the pre-core orbit, the linear centered terms vanish.  The
quadratic terms are the joint-success kernel on plus portions minus
$1/p^2$ times the unrestricted kernel on the same pre-link
portions, with the signed endpoint kernels just described.  Both
terms carry the \emph{product} of their after-mask probabilities.
After summing locations, orbit-averaged joint successes can
equivalently be counted as actual simultaneous successes.

Whenever an endpoint pair contributes to a broad signed kernel,
the two corresponding plus positions are within $O(r_{xy}^+)$,
by the cell-span bound.  The total link fractions emanating from
one plus label are at most its mass.  Therefore all the absolute
$1/p^2$ terms and all joint-success terms outside the tightest
range are bounded by the first two parts of
Proposition~\ref{tab:sparse-input}.

\emph{Tight cancellation.}\quad
At distance at most
$r_{xy}^+\exp(-\exp(\sigma_gY))$ between the plus positions,
all endpoint distances, including those involving matched minus
positions, qualify for the broad kernel.  More explicitly, choose
$k_0$ small enough that the constant in \eqref{tab:cell-span}
times $e^{-E_gB_i}$ is at most $1/10$ for every retained centre.
Each linked displacement is then at most one tenth of its own
width.  All four endpoint distances are at most
$(1+1/10+1/10)r_{xy}^+<10r_{xy}^+$.  Matching preserves each
endpoint's width, so their four kernel values, including the
factor $r_{xy}^-$, are identical.  The signed contribution
of a transfer consequently cancels.  Only products of genuinely
unpaired plus arrivals remain in this range.

\emph{One-width self cost.}\quad
First fix one exact width $r$.  Consider the self cost consisting
of products of aligned unpaired plus arrivals within the relaxed
range $r\exp(-\tfrac12\exp(\sigma_gY))$, multiplied by $r$.
Pairs in different old classes are paid either by exceptions in
Proposition~\ref{tab:sparse-input} or by the coarsening increment
$K_p'$.

For the within-class pairs, let $P_A$ be the successful arriving
plus mass in an old class $A$ that is either unpaired or matched
to a minus portion outside $A$.  Let $N_A$ be the mass removed
from minus portions in $A$ by successful plus portions outside
$A$.  These are weighted masses, including their after-mask
probabilities.  Internal transfers cancel, and unpaired minus
weights have no centered fluctuation.  The old class fluctuation
is therefore the centered version of $P_A-N_A$.  Orbit
expectation gives
\begin{align}
 r\sum_A\mathbb E P_A^2
 &\le J_p+
 r\sum_A\left\{\bigl(\mathbb E(P_A-N_A)\bigr)^2
                    +2\mathbb E(P_AN_A)\right\}.
 \label{tab:self-cost}
\end{align}
To check the inequality, expand
$P_A^2=(P_A-N_A)^2+2P_AN_A-N_A^2$ and discard the last,
nonpositive term.  The variance part is included in $J_p$.

\emph{Capacity bounds.}\quad
Here are explicit capacity bounds for both charges.  Index plus
sources by $\xi$, with old class $B_\xi$, success indicator
$S_\xi$, and mass $c_\xi$ that it would carry upon success
(its residual-scaled raw weight times its after-mask probability).
Let $q_{\xi A}$ be the weighted mass of its links to minus
portions in $A\ne B_\xi$, and put $q_{\xi B_\xi}=0$.  Then
\[
 \sum_Aq_{\xi A}\le c_\xi,\qquad
 P_A\le\sum_{B_\xi=A}c_\xi S_\xi,
 \qquad N_A=\sum_\xi q_{\xi A}S_\xi.
\]
Every link uses the same after-mask at its two ends; the
inequalities therefore remain true when $A$ contains different
masks.  With
$a_{\xi A}=c_\xi\mathbf1_{\{B_\xi=A\}}+q_{\xi A}$,
we have $\sum_Aa_{\xi A}\le2c_\xi$.  Consequently,
\begin{align*}
 \sum_A\bigl|\mathbb E(P_A-N_A)\bigr|^2
 &\le p^{-2}\sum_{\xi,\eta}\sum_Aa_{\xi A}a_{\eta A}\\
 &\le 4p^{-2}\sum_{\substack{\xi,\eta\\
                  \operatorname{dist}(\xi,\eta)=O(r)}}c_\xi c_\eta,\\
 \sum_AP_AN_A
 &\le\sum_{B_\xi\ne B_\eta}
       c_\xi q_{\eta B_\xi}S_\xi S_\eta\\
 &\le\sum_{\substack{B_\xi\ne B_\eta\\
                  \operatorname{dist}(\xi,\eta)=O(r)}}
       c_\xi c_\eta S_\xi S_\eta.
\end{align*}
The distance restrictions are justified before enlarging the
sums: a nonzero term in a common class $A$ has both source
locations in its cell, either directly or through a link.
The first bound is paid by the before-plus counts times
$1/p^2$.  In the second, each ordered source pair occurs only
once, regardless of how many destination classes receive parts
of one source.  The far or exceptional pairs are paid by
Proposition~\ref{tab:sparse-input}.  Each remaining pair is in
the relaxed range and joins two old classes, so its charge is
at most the corresponding product of retained post-class masses
included in $K_p'$.  Summing \eqref{tab:self-cost} over widths,
components, and primes and applying
Lemma~\ref{tab:variance-budget} bounds the total relaxed self cost.

\emph{Unequal widths.}\quad
Finally take unequal widths.  Bin the actual post-projection
locations in bins of length $\delta=r_{\max}e^{-E_Y}$, where
$E_Y=\exp(\sigma_gY)$.  By \eqref{eq:component-path},
\[
 \log(r_{\max}/r_{\min})+\log3
 \le \exp(h_*Y/10)+\log3<\tfrac12E_Y
\]
for all permitted heights, so $3\delta$ is below the smaller
width's relaxed self scale.  Let $X_{i,b}$ be the unpaired
arriving mass of width $r_i$ in bin $b$, and let $\mathcal S_i$
be its full relaxed self cost at this step.  A tight pair is
in the same or adjacent bins, whence
\[
 r_{\min}\sum_{|b-c|\le1}X_{i,b}X_{j,c}
 \le3r_{\min}
       \left(\sum_bX_{i,b}^2\sum_cX_{j,c}^2\right)^{1/2}
 \le3\sqrt{\frac{r_{\min}}{r_{\max}}}
       \sqrt{\mathcal S_i\mathcal S_j}.
\]
First apply Cauchy--Schwarz over all steps and components for
each fixed pair of widths.  Then sum the resulting kernel
$2^{-|i-j|/2}$ over dyadic width indices; its summable convolution
bounds the total by a constant times the total self cost.  No
accounting class has to mix widths.  This bounds every sparse
increment of the broad energy; its birth and join injections are
bounded by Proposition~\ref{tab:first-join-input}.
\end{proof}

\subsection{Virtual weights and first-mass losses}

The energy argument supplies a uniform second moment.  To obtain
a positive integral on a fixed set, we also need to understand the
spatial distribution of first mass.  Adaptive weights need not be
fixed residue exclusions.  Their exact orbit marginals provide a
replacement for that simpler description.

\begin{proposition}[Virtual product marginals]
\label{tab:virtual-marginals}
For each retained good row, there is a nonnegative virtual weight
on its true numerator grid with the following properties.
\begin{enumerate}
\item It dominates the actual final weight, and its first mass
lies between $c_{\rm hard}\phi(v)r(v)$ and $\phi(v)r(v)$.
\item Its total excess over the actual final first mass is exactly
accounted for by the first-moment discard losses in the projected
tables, when the virtual weight omits the soft and additional
static tests, cap clipping, and partition discards.
\item At fixed static anchor choices, every subset of the
actual hard and adaptive fields has mean equal to the product
of its local means.  An actual hard field has mean $1-b_p/p$
when it forbids $b_p\le2$ distinct residues, or mean $R_p$
when it is a prescribed scalar; adaptive full fields have mean
$R_p$.  For upper counts one may retain just one hard birth
exclusion at each applicable prime, or the scalar $R_p$ when
there is no exclusion.  Any subset of the resulting hard and
adaptive fields has the product of its ordinary local means.
\item Fields through a prime $p$, restricted to a later projection
of denominator $v'$, have a joint period dividing $v'$ of size
at most $\exp(O(p))$.
\end{enumerate}
Domination in the first item is after the static anchor average;
it is not required separately for each auxiliary anchor choice.
\end{proposition}

\begin{proof}
We first construct the fields and prove domination and the first-mass
bounds.  We then verify subset means, identify discard losses by
telescoping, and control the periods.

\emph{Field construction and domination.}\quad
Fix a row and a location on its true grid.  At each full adaptive
prime $p$, take the actual fractional multiplier of its minus
ancestor at the $p$-update as the normalized field at that
location.  If the entry or plan is unavailable there, extend the
field by $R_p$.  Unavailability is pre-$p$-equivariant.
For clarity, if its positive raw prefix weight is $w_x$, its
link masses are $m_{x\xi}$, and
$u_x=w_x-\sum_\xi m_{x\xi}$ is unpaired, its multiplier is
\[
 M_x=1-\sum_\xi\frac{m_{x\xi}}{w_x}S_\xi
          -\frac{u_x}{p w_x}.
\]
The coefficients and the positive-prefix condition transport
invariantly on the pre-$p$ orbit.  Since each $S_\xi$ has mean
$1/p$, the field has orbit mean $R_p$; the extension has the
same mean on unavailable or zero-prefix orbits.  These are
translations by $1/p^\nu$, where $p^\nu\parallel L_i$.
All earlier fields are invariant along this orbit.

Because this row is full at $p$, its later grid is preserved by
the same translations.  In particular, restricting to later
omission subarrays does not destroy the orbit: the increment of
the earlier numerator is divisible by their prime factors,
which are coprime to $p$.  At an adaptive omission, the later
grid itself implements the finer projection; the only remaining
field at this prime is the scalar $R_p$ if its residual power is
positive.

Multiply these adaptive fields by the full hard birth weight,
omitting soft tests, the additional static discard, clipping,
and partition discards.  To handle hard anchor averaging, first
fix all static anchor choices but use the \emph{same} adaptive
multiplier functions computed from the actual averaged tables.
Every such fixed earlier hard field is invariant along each
future prime orbit.  Independence of the anchor choices from
the adaptive plans is therefore unnecessary.  Average the
resulting virtual products over the anchor choices at the end.

For the actual averaged hard weight, the actual final weight is
its product with these multipliers and the intervening discard
fractions.  All discard fractions lie in $[0,1]$, so the virtual
weight dominates it.  Reverse orbit averaging removes the
adaptive fields one by one.  CRT then gives the hard birth
means, bounded between $c_{\rm hard}$ and $1$ times the ordinary
totient density.  The true-grid omission proportions supply
exactly the missing mask factors, as in
\eqref{tab:first-currency}.  This proves the first item.

\emph{Subset means.}\quad
The same reverse averaging works for any selected subset of
the actual fields: remove the selected adaptive fields in
reverse order, retaining all selected earlier hard fields,
then use CRT on the hard fields.  Their residue colours may
depend on the fixed anchors, but their means are $1-b_p/p$
with $b_p\le2$, or $R_p$ for a scalar.  All these means are
positive: two active distinct exclusions occur only at large
primes by the model-separation estimate.  No lower bound on
their full product is needed for this exact subset identity.
For an upper version choose, within each fixed static
anchor prescription, one of the exclusions at a prime with an
active exclusion, independently of the partner in a pair
count.  Its mean is $R_p$; if none is active keep the prescribed
scalar $R_p$.  Dropping other exclusions and filters gives a
pointwise upper bound.  Any retained subset of these fields and
of the adaptive fields again has product mean, proving the
third item.

\emph{Discard telescoping.}\quad
For the second item, pull each true-grid label back through
its same-location ancestors and order its finite operations.
Let $f_j$ denote the nondiscard multiplier at operation $j$,
and $d_j\in[0,1]$ its retained discard fraction, inserting
identity factors when only one operation occurs.  Take the
anchor-averaged hard birth weight as the initial factor.
The exact finite identity is
\[
 \prod_{j=1}^n f_j-\prod_{j=1}^n(f_jd_j)
 =\sum_{j=1}^n
    \left(\prod_{k<j}f_kd_k\right)f_j(1-d_j)
    \left(\prod_{k>j}f_k\right).
\]
Thus every term has the actual prefix and the virtual suffix,
even when adaptive plans depend on the actual preceding
tables.  Future full multipliers eliminate in
reverse because the discard decision and its prefix are
invariant under the corresponding unprocessed-prime
translations.  Future omitted coordinates restrict to their
true numerator subarrays and contribute their residual
scalars.  More explicitly, for a remaining mask $a\mid I$,
averaging the future full fields on that subarray leaves the
factor $R(I/a)$.  The current prefix-weighted discard array
meets the subarray in exactly the fraction $1/a$, by its
still-full unprocessed core.  The factor recovered at the
discard time is thus
\begin{equation}
 \frac{R(I/a)}a=\pi_I(a).
 \label{tab:discard-currency}
\end{equation}
This is precisely its present first-mass currency.  Summing
these telescoping terms proves the assertion.  Soft-test losses
may accordingly be compared to hard birth means; future fields
restore those same means.

\emph{Period control.}\quad
Finally an adaptive field through $p$ is invariant under the
post-$p$ core translations.  Intersect this translation group
with a later grid of denominator $v'$.  Its period in the later
numerator coordinate divides $v'$ and is at most
\[
 \frac{v'}{(v',\mathrm{core}_p)}
 \le\frac{L_i}{\mathrm{core}_p}\le\exp(O(p)),
\]
by \eqref{eq:core-loss}.  Earlier field cores are divisible
by the post-$p$ core, so this also gives a common period for
all fields through $p$.  The constant is fixed by the hierarchy
and independent of the block.  Static hard fields use only
residue colours modulo their respective primes, so their joint
modulus divides $\prod_{l\le p}l\le\exp(O(p))$.  The potentially
large rational model numerator and denominator specify the
colour, not its modulus; if the denominator is divisible by
the prime, that exclusion is inactive.  Soft and hole tests
are omitted from the virtual product.  Thus their larger
allowed moduli and the hole cutoff $\lambda_0$ do not affect
this period bound.  This proves
the fourth item.
\end{proof}

We will also use the preceding argument before the last projection.
The following version makes explicit the interface needed to count an
existing table against a fixed spatial query in
Sections~\ref{sec:partitions} and~\ref{sec:joins}.

\begin{corollary}[Upper fields on an intermediate projection]
\label{tab:intermediate-fields}
Fix an admissible realized table history and a row whose current
projected grid is $(z+a\gamma)/(L_i/d)$, where $I$ is the remaining
prime-power part and $a\mid I$. On this grid the actual raw weight
is bounded, after averaging static anchors, by a product of fields
at the processed primes. Each nonconstant field lies in $[0,1]$.
At a full adaptive prime use the multiplier computed from the actual
matching, extended by $R_p$ where its prefix is unavailable; at a
birth prime retain one active hard exclusion, or its prescribed scalar
if none is active. An adaptive omission contributes only its residual
scalar when a positive power remains.

For every fixed choice of the static anchors, every subset of these
fields has mean equal to the product of its ordinary local means.
In particular the full upper product has mean
\[
 R_{\mathrm{done}}(L_i/d)
 :=R\bigl((L_i/I)/d\bigr).
\]
The joint numerator period of any retained fields through a prime $p$
divides the current projected denominator and is at most $\exp(O(p))$,
with the same fixed-parameter uniformity as in
Proposition~\ref{tab:virtual-marginals}.
\end{corollary}

\begin{proof}
Stop the field construction in
Proposition~\ref{tab:virtual-marginals} at the present projection.
Omit all discard fractions and keep the adaptive multiplier functions
from the realized tables. Choosing one hard exclusion can only enlarge
the fixed-anchor birth weight; average these enlarged products after
multiplying by the same adaptive fields. This gives the claimed
pointwise upper bound on the actual raw weight.

Reverse orbit averaging applies to every selected adaptive field
already processed. Later omissions, now only those before the present
time, use other primes and preserve its full-prime translation orbit.
The remaining static fields have their ordinary means by CRT. A
processed adaptive omission leaves precisely the residual scalar prescribed by
the prime update. Thus every subset has its product mean, and the
product over all primes is $R((L_i/I)/d)$. Finally intersect the
historical post-$p$ core translations with the present grid, exactly
as in the period-control argument above. Removing some fields does
not alter the remaining field functions or enlarge their joint period.
\end{proof}

\begin{lemma}[Equidistribution of virtual first mass]
\label{tab:equidistribution}
On every fixed nondegenerate spatial interval, the virtual
weighted mass is the interval's length times the total virtual
row mass, with relative error tending to zero uniformly as the
denominators tend to infinity.  The same statement holds for
their weighted trains
of interval indicators of half-width $r(v)/2$.
\end{lemma}

\begin{proof}
On the true grid of denominator $v$, first retain only fields
at primes $p\le c_0\log v$, with $c_0>0$ sufficiently small
and fixed.  Proposition~\ref{tab:virtual-marginals} gives a
joint period dividing $v$ and at most $v^{1/2}$.  For a
nonnegative periodic array, complete periods inside a fixed
interval contribute exactly its period mean; the two possible
incomplete periods give relative error $O(v^{-1/2})$.  This
constant may depend on the fixed interval's positive length;
the bound does not require a lower bound on the mean of one period.

Deleting the remaining fields changes the total mean by at
most a relative constant times
\[
 \sum_{\substack{p\mid v\\p>c_0\log v}}\frac1p=o(1).
\]
Indeed there are at most
$\log v/\log(c_0\log v)$ such prime divisors, so the sum is
$O(1/\log\log v)$.  To see the relative comparison explicitly,
fix anchors and denote the full and truncated means by $M$ and
$M_0$.  The actual-field subset identity of
Proposition~\ref{tab:virtual-marginals} gives
\[
 \frac{M_0-M}{M}
 =\prod_{\substack{p\mid v\\p>c_0\log v}}\mu_p^{-1}-1
 \le \exp\left(C\sum_{\substack{p\mid v\\p>c_0\log v}}\frac1p\right)-1,
\]
where $\mu_p\ge1-2/p$ is the relevant local mean.  All low
means cancel from this ratio, even if their product is very
small.  This uniform fixed-anchor estimate survives anchor
averaging.  Deleting factors in $[0,1]$ increases the array, so
this is an $L^1$ comparison and proves equidistribution for the
full virtual array.  Work first with fixed anchor choices and
then average; all estimates are uniform in those choices.

Since $r(v)\to0$, integrals of the weighted interval trains are
squeezed between the corresponding point counts in slightly
enlarged and slightly shortened fixed intervals.  This gives
the final assertion.
\end{proof}

\subsection{Completion of the conditional reduction}

\begin{proposition}[Finite-block reduction]
\label{tab:conditional-reduction}
Assume the selections and birth prescriptions of
Sections~\ref{sec:centres} and~\ref{sec:prescriptions}, and
Propositions~\ref{tab:first-join-input}--\ref{tab:partition-input}.
Then no measurable set of positive measure can avoid the
original approximation intervals in a whole tail.  Consequently
these inputs imply Theorem~\ref{thm:main}.
\end{proposition}

\begin{proof}
Let the proposed avoided set $E$ have measure $\mu>0$.  The row
selection and full-mask purges retain a fixed fraction $c_*>0$
of the block mass.  Indeed, after the raw-row case has been
excluded, the nonraw rows have a fixed positive share; the
negligible purges and selection of one of finitely many types
and colours preserve a fixed share.  The constant
$c_{\rm hard}>0$ is fixed by the hard prescriptions before
the later loss budgets are chosen.  Let $F^{\rm vir}$ be the sum of the virtual
weighted indicators with half-width $r(v)/2$, and let $F$ be
the corresponding realized sum on the true grids.  An indicator
has integral $r(v)$, so its mean integral agrees with the first
mass convention.  Proposition~\ref{tab:virtual-marginals} gives
\[
 M_-\le\int F^{\rm vir}\le M_+,\qquad
 M_-=c_{\rm hard}c_*w,\quad M_+=2w.
\]
These bounds are uniform in the block and every realized
partition history; we will not need to assume that its virtual
mass is deterministic.

The soft loss must be compared with the actual virtual first mass,
rather than with its uniform lower bound. Put
\[
 V=\mathbb E\int F^{\rm vir},\qquad M_-\le V\le M_+.
\]
By Lemma~\ref{pre:soft-loss} and the discard telescoping identity,
its contribution to the expected final loss is at most $\eta_s V$,
where $\eta_s=C\epsilon_s/h_s\le\mu/20$. This is the choice of
$\epsilon_s$ already made in the hierarchy; it does not require
$\epsilon_s$ to be selected after the resulting $c_{\rm hard}$.

With $c_{\rm hard}$ now fixed, make the additional static and partition
losses small using Propositions~\ref{tab:first-join-input}
and~\ref{tab:partition-input}. The clipping loss is at most
\[
 \frac1{N_0}\{O(W)+O_{k_0}(W^{1+c})\},
\]
by \eqref{tab:clip-child} and \eqref{tab:narrow-input}.
Choose the fixed cap only after those comparison constants.
Since $W\le2w$, the later loss budgets and the cap can be chosen so
that the combined non-soft loss is at most $\mu M_-/20$. Thus
\[
 \mathbb E\|F^{\rm vir}-F\|_1
 \le \eta_s V+\mu M_-/20.
\]
If no jitters are used, omit the expectation. The telescoping identity
\eqref{tab:discard-currency} ensures that these are exactly the
relevant final losses, including the soft loss measured against the
hard birth means.

Proposition~\ref{tab:energy-bound} gives a uniform bound on
$\|F\|_2$.  Indeed, the intersection length of two half-width
intervals is at most $r_{xy}^-$ and vanishes unless their
centres are at distance at most $r_{xy}^+$.  It is therefore
bounded by the final broad energy.  If jitters are used, the
same bound holds for $\mathbb E F$ by Jensen's inequality.
Fix a uniform constant $C_2$ with $\|\mathbb E F\|_2\le C_2$.

All constants, including $N_0,C_2,M_-$, are now fixed.  Choose
a fixed finite union of intervals $U$ with
\[
 |U\mathbin{\triangle}E|
 \le \min\left\{\frac\mu2,\,
           \left(\frac{\mu M_-}{4C_2}\right)^2\right\}.
\]
In particular $|U|\ge\mu/2$.  Uniform relative equidistribution
from Lemma~\ref{tab:equidistribution}, applied only after this
choice, gives on every sufficiently late block
\[
 \mathbb E\int_U F^{\rm vir}
 \ge (|U|-\mu/10)V\ge .4\mu V.
\]
Subtracting the loss bound gives
\[
 \int_U\mathbb E F
 \ge (.4\mu-\eta_s)V-\mu M_-/20
 \ge .3\mu M_-.
\]  However, every realized
summand is supported inside an original approximation interval,
so $\int_E\mathbb E F=0$.  Cauchy--Schwarz then gives
\[
 \int_U\mathbb E F
 \le C_2|U\mathbin{\triangle}E|^{1/2}\le .25\mu M_-,
\]
a contradiction.  This completes the
conditional reduction.
\end{proof}

To finish the argument, it remains to establish
Propositions~\ref{tab:first-join-input}, \ref{tab:sparse-input},
and~\ref{tab:partition-input}; Section~\ref{sec:joins} assembles
these inputs in the required chronological order.
